
\section{Tote Zone I}
\label{section:tote_zone_1}


\begin{frame}
\begin{columns}
    \begin{column}{0.58\textwidth}
    
\begin{figure}
    \DARCimage{1.0\linewidth}{741include}
    \caption{\scriptsize Die Tote Zone, die für die Bodenwelle zu nah und für die Raumwelle zu weit weg ist.}
    \label{n_ionosphaere_tote_zone}
\end{figure}


    \end{column}
   \begin{column}{0.38\textwidth}
       \begin{itemize}
        \item Bereich, wo die \emph{Bodenwelle nicht mehr} hin gelangt
        \item Und die \emph{Raumwelle noch nicht} hingelangt
        \item Abhängig vom Reflexionswinkel der Raumwelle
        \item Funkstationen in der Toten Zone können mich nicht hören
  \end{itemize}

   \end{column}
\end{columns}

\end{frame}

\begin{frame}
\only<1>{
\begin{QQuestion}{BE106}{Eine Frequenz auf einem höheren Kurzwellenband erscheint zunächst frei, stellt sich aber anschließend als besetzt heraus. Was ist die häufigste Ursache dafür?}{Die auf dieser Frequenz sendende Station wurde durch den Mögel-Dellinger-Effekt kurzfristig unterbrochen.}
{Eine Station auf dieser Frequenz verwendet das andere Seitenband.}
{Für die auf dieser Frequenz sendenden Stationen sind die Ausbreitungsbedingungen zu schlecht.}
{Eine auf dieser Frequenz sendende Station liegt innerhalb der toten Zone und konnte daher von mir nicht gehört werden.}
\end{QQuestion}

}
\only<2>{
\begin{QQuestion}{BE106}{Eine Frequenz auf einem höheren Kurzwellenband erscheint zunächst frei, stellt sich aber anschließend als besetzt heraus. Was ist die häufigste Ursache dafür?}{Die auf dieser Frequenz sendende Station wurde durch den Mögel-Dellinger-Effekt kurzfristig unterbrochen.}
{Eine Station auf dieser Frequenz verwendet das andere Seitenband.}
{Für die auf dieser Frequenz sendenden Stationen sind die Ausbreitungsbedingungen zu schlecht.}
{\textbf{\textcolor{DARCgreen}{Eine auf dieser Frequenz sendende Station liegt innerhalb der toten Zone und konnte daher von mir nicht gehört werden.}}}
\end{QQuestion}

}
\end{frame}%ENDCONTENT


\begin{frame}
\only<1>{
\begin{QQuestion}{EH201}{Unter der \glqq Toten Zone\grqq{} wird der Bereich verstanden,~...}{der durch die Bodenwelle erreicht wird und für die Raumwelle nicht zugänglich ist.}
{der durch die Bodenwelle überdeckt wird, so dass schwächere DX-Stationen zugedeckt werden.}
{der durch die Bodenwelle nicht mehr erreicht wird und durch die Raumwelle noch nicht erreicht wird.}
{der durch die Überlagerung der Bodenwelle mit der Raumwelle in einer Zone der gegenseitigen Auslöschung liegt.}
\end{QQuestion}

}
\only<2>{
\begin{QQuestion}{EH201}{Unter der \glqq Toten Zone\grqq{} wird der Bereich verstanden,~...}{der durch die Bodenwelle erreicht wird und für die Raumwelle nicht zugänglich ist.}
{der durch die Bodenwelle überdeckt wird, so dass schwächere DX-Stationen zugedeckt werden.}
{\textbf{\textcolor{DARCgreen}{der durch die Bodenwelle nicht mehr erreicht wird und durch die Raumwelle noch nicht erreicht wird.}}}
{der durch die Überlagerung der Bodenwelle mit der Raumwelle in einer Zone der gegenseitigen Auslöschung liegt.}
\end{QQuestion}

}
\end{frame}%ENDCONTENT
